
\documentclass{article}
\usepackage{colortbl}
\usepackage{makecell}
\usepackage{multirow}
\usepackage{supertabular}

\begin{document}

\newcounter{utterance}

\centering \large Interaction Transcript for game `bbh', experiment `logical\_deduction\_three\_objects', episode 53 with reflexion{-}Qwen3.5{-}4B.
\vspace{24pt}

{ \footnotesize  \setcounter{utterance}{1}
\setlength{\tabcolsep}{0pt}
\begin{supertabular}{c@{$\;$}|p{.15\linewidth}@{}p{.15\linewidth}p{.15\linewidth}p{.15\linewidth}p{.15\linewidth}p{.15\linewidth}}
    \# & \multicolumn{2}{c}{Player} && \multicolumn{2}{c}{Game Master} \\
    \hline

    \theutterance \stepcounter{utterance}  
    & & & \multicolumn{4}{p{0.6\linewidth}}{
        \cellcolor[rgb]{0.9,0.9,0.9}{
            \makecell[{{p{\linewidth}}}]{
                \texttt{\tiny{[P1$\langle$GM]}}
                \texttt{You have attempted this same task before and did not succeed. These are the plans you wrote after those attempts.} \\
\texttt{Trial \#1: New plan: I failed to correctly interpret the phrase "second{-}newest" as the middle position in a three{-}item sequence, leading to a logical contradiction when trying to place the tractor. Next time, I will explicitly map the relative terms to fixed positions (1st, 2nd, 3rd) immediately upon reading the prompt. Specifically, I will treat "second{-}newest" as the 2nd position, then use the "older than" constraint to fill the remaining spots, ensuring I verify that the final order satisfies all conditions before selecting the option.} \\
\\ 
\\ 
\texttt{A logical deduction task which requires deducing the order of a sequence of objects.} \\
\\ 
\texttt{Q: The following paragraphs each describe a set of three objects arranged in a fixed order. The statements are logically consistent within each paragraph. In a golf tournament, there were three golfers: Amy, Eli, and Eve. Eve finished above Amy. Eli finished below Amy.} \\
\texttt{Options:} \\
\texttt{(A) Amy finished last} \\
\texttt{(B) Eli finished last} \\
\texttt{(C) Eve finished last} \\
\texttt{A: (B)} \\
\\ 
\texttt{Q: The following paragraphs each describe a set of three objects arranged in a fixed order. The statements are logically consistent within each paragraph. On a shelf, there are three books: a white book, a green book, and an orange book. The green book is to the right of the white book. The orange book is the rightmost.} \\
\texttt{Options:} \\
\texttt{(A) The white book is the leftmost} \\
\texttt{(B) The green book is the leftmost} \\
\texttt{(C) The orange book is the leftmost} \\
\texttt{A: (A)} \\
\\ 
\texttt{Q: The following paragraphs each describe a set of three objects arranged in a fixed order. The statements are logically consistent within each paragraph. On a shelf, there are three books: a red book, a gray book, and a white book. The white book is to the left of the gray book. The red book is the second from the left.} \\
\texttt{Options:} \\
\texttt{(A) The red book is the leftmost} \\
\texttt{(B) The gray book is the leftmost} \\
\texttt{(C) The white book is the leftmost} \\
\texttt{A: (C)} \\
\\ 
\texttt{Q: The following paragraphs each describe a set of three objects arranged in a fixed order. The statements are logically consistent within each paragraph. In an antique car show, there are three vehicles: a motorcyle, a minivan, and a tractor. The minivan is older than the tractor. The minivan is the second{-}newest.} \\
\texttt{Options:} \\
\texttt{(A) The motorcyle is the oldest} \\
\texttt{(B) The minivan is the oldest} \\
\texttt{(C) The tractor is the oldest} \\
\texttt{A:} \\
            }
        }
    }
     \\ \\

    \theutterance \stepcounter{utterance}  
    & \multicolumn{4}{p{0.6\linewidth}}{
        \cellcolor[rgb]{0.561,0.737,0.561}{
            \makecell[{{p{\linewidth}}}]{
                \texttt{\tiny{[P1$\rangle$GM]}}
                \texttt{Based on the plan to explicitly map relative terms to fixed positions:} \\
\\ 
\texttt{1.  **Analyze the sequence**: There are three vehicles: motorcyle, minivan, and tractor. Let's assign positions 1 (oldest), 2 (middle), and 3 (newest).} \\
\texttt{2.  **Map "second{-}newest"**: In a sequence of three, "second{-}newest" corresponds to the middle position (Position 2). Therefore, the **minivan is in the middle**.} \\
\texttt{3.  **Apply the "older than" constraint**: The statement says "The minivan is older than the tractor." Since the minivan is in the middle (Position 2), the tractor must be older than it, placing the **tractor in Position 1 (the oldest)**.} \\
\texttt{4.  **Deduce the remaining item**: This leaves the motorcyle for Position 3 (the newest).} \\
\texttt{5.  **Final Order**:} \\
\texttt{    *   1st (Oldest): Tractor} \\
\texttt{    *   2nd (Middle): Minivan} \\
\texttt{    *   3rd (Newest): Motorcyle} \\
\texttt{6.  **Evaluate Options**:} \\
\texttt{    *   (A) The motorcyle is the oldest {-}> False (it is the newest).} \\
\texttt{    *   (B) The minivan is the oldest {-}> False (it is the middle).} \\
\texttt{    *   (C) The tractor is the oldest {-}> True.} \\
\\ 
\texttt{A: (C)} \\
            }
        }
    }
    & & \\ \\

    \theutterance \stepcounter{utterance}  
    & & & \multicolumn{2}{p{0.3\linewidth}}{
        \cellcolor[rgb]{0.9,0.9,0.9}{
            \makecell[{{p{\linewidth}}}]{
                \texttt{\tiny{[GM$|$GM]}}
                \texttt{Target: (A)} \\
            }
        }
    }
    & & \\ \\

    \theutterance \stepcounter{utterance}  
    & & & \multicolumn{2}{p{0.3\linewidth}}{
        \cellcolor[rgb]{0.9,0.9,0.9}{
            \makecell[{{p{\linewidth}}}]{
                \texttt{\tiny{[GM$|$GM]}}
                \texttt{game\_result = LOSE} \\
            }
        }
    }
    & & \\ \\

    \theutterance \stepcounter{utterance}  
    & & \multicolumn{4}{p{0.6\linewidth}}{
        \cellcolor[rgb]{0.95,0.95,0.95}{
            \makecell[{{p{\linewidth}}}]{
                \texttt{\tiny{[P1$\leftrightarrow$P]}}
                \texttt{Post{-}episode Reflexion (not sent to GM):} \\
\texttt{New plan: I failed to correctly interpret the phrase "second{-}newest" as the middle position in a three{-}item sequence, leading to a logical contradiction when trying to place the tractor. Next time, I will explicitly map the relative terms to fixed positions (1st, 2nd, 3rd) immediately upon reading the prompt. Specifically, I will treat "second{-}newest" as the 2nd position, then use the "older than" constraint to fill the remaining spots, ensuring I verify that the final order satisfies all conditions before selecting the option.} \\
            }
        }
    }
    & \\ \\

\end{supertabular}
}

\end{document}
